\documentclass[12pt,a4paper]{report}
\usepackage[margin=1in]{geometry}
\usepackage{setspace}
\usepackage{titlesec}
\usepackage{hyperref}
\usepackage{enumitem}
\usepackage[T1]{fontenc}
\usepackage[utf8]{inputenc}
\onehalfspacing
\begin{document}

\title{FormFix AI\\Real-Time Exercise Pose Detection, Rep Counting, and Posture Correction System}
\author{Syed Ahmed}
\date{}
\maketitle

\tableofcontents

\chapter{Introduction}
\section{1.1 Introduction}
Artificial intelligence has rapidly transformed many application areas including healthcare, surveillance, education, and fitness. In the domain of personal health and training, AI systems are increasingly being used to automate exercise analysis, estimate body posture, detect movement quality, and guide users without constant dependence on a human trainer. At the same time, the popularity of home workouts and digital fitness programs has increased significantly, creating a strong need for low-cost intelligent workout assistance systems. Many people attempt exercises from online videos or mobile applications, but they often perform these movements with incorrect posture, incomplete range of motion, poor joint alignment, or unstable form. Such mistakes may reduce exercise effectiveness and, in some cases, increase the risk of injury.

Traditional fitness applications usually focus on workout scheduling, calorie tracking, or static exercise descriptions. Although these features are useful, they do not actively observe the user and therefore cannot verify whether the user is actually performing the movement correctly. Human personal trainers can provide this level of guidance, but personal training is expensive, not always accessible, and difficult to scale for everyday use. This creates a gap between instructional content and actual exercise quality. An intelligent real-time fitness system that can watch the user, understand the movement being performed, count repetitions, detect form errors, and provide corrective feedback can help bridge this gap.

The proposed project, FormFix AI, is developed to address this need. It is an AI-based virtual fitness assistant that observes a user through a webcam, extracts skeletal landmarks using pose estimation, converts these landmarks into angle-based motion features, classifies the exercise being performed using sequence models, counts repetitions using stage transitions, and evaluates posture using biomechanical rules. The system also includes voice feedback and detailed end-of-exercise reporting so that the user not only receives correction in the moment but also understands what was done correctly and what still needs improvement. This combination makes the project both technically meaningful and practically useful.

\section{1.2 Objective}
The main objective of this project is to design and implement an intelligent exercise monitoring system capable of recognizing exercises in real time, counting repetitions accurately, detecting posture deviations, and delivering corrective feedback in a human-friendly manner. The goal is not only to identify which exercise is being performed, but also to understand the quality of the performance and guide the user toward safer and more effective movement.

The specific objectives of the project are as follows:

\noindent\textbullet\ To detect human body landmarks and generate a live skeleton representation using computer vision.
\noindent\textbullet\ To classify supported exercises such as bicep curls, push-ups, pull-ups, barbell squats, and lateral raises using temporal deep learning models.
\noindent\textbullet\ To count repetitions accurately by tracking angle transitions and motion stages.
\noindent\textbullet\ To analyze body posture in real time and identify common mistakes based on exercise-specific rules.
\noindent\textbullet\ To deliver real-time visual and voice-based corrective feedback.
\noindent\textbullet\ To generate a detailed exercise report showing performance quality, perfect repetitions, corrected repetitions, and form issues.
\noindent\textbullet\ To maintain user progress through a SQL-based backend that stores completed, skipped, and in-progress exercises.

\section{1.3 Scope of the Project}
The scope of the project is limited to a webcam-based real-time exercise trainer for selected strength and bodyweight exercises. The current version of the system is designed for indoor use where a user performs exercises in front of a single camera. The system primarily focuses on exercises whose motion can be reasonably represented using skeletal landmarks and angle-based analysis. These include bicep curls, push-ups, pull-ups, barbell squats, and lateral raises. The project does not attempt to cover all gym exercises, especially those with heavy machine occlusion, extreme side-angle dependency, or complex object interactions that are difficult to analyze from a single webcam.

From the software perspective, the scope includes a training pipeline for sequence models, a real-time inference engine, a frontend workout interface, a backend API, a database for progress tracking, and an optional local large language model based coaching layer for human-like feedback. The project is intended as a deployable academic and prototype fitness assistant rather than as a certified medical or rehabilitation device. Therefore, although it can support safe general exercise guidance, it should not be treated as a substitute for professional diagnosis, medical supervision, or advanced sports biomechanics testing.

\section{1.4 Problem Statement}
Many users perform exercise routines without proper supervision. As a result, they may unknowingly use poor form, incomplete movement range, incorrect posture, or unstable control, which can reduce the effectiveness of the exercise and increase the possibility of fatigue or injury. Existing fitness applications generally provide workout plans and timers but often fail to monitor the actual execution quality of the exercise in real time. On the other hand, direct access to professional personal trainers is expensive, limited by time and location, and not always feasible for students or home users.

The problem addressed in this project is the lack of an affordable, real-time, intelligent workout monitoring system that can identify the exercise being performed, count repetitions accurately, detect posture mistakes, guide the user with corrective feedback, and store progress for long-term improvement. The challenge lies in achieving this using consumer hardware with sufficient speed, interpretability, and practical accuracy.

\section{1.5 Motivation}
The motivation behind this project comes from the growing importance of fitness, the popularity of home workouts, and the increasing role of AI in human-centered systems. Many users want trainer-like guidance but do not have access to one every day. At the same time, modern computer vision and temporal deep learning models make it possible to build systems that understand body movement from simple video input. This project aims to bring these technologies together in a useful and accessible form.

Another strong motivation is to build a system that goes beyond simple exercise classification. Recognizing that a user is doing a bicep curl is useful, but understanding whether the elbow is stable, whether the curl reaches full contraction, whether the movement is balanced, and whether the user is maintaining proper control is far more valuable. By combining pose estimation, sequence learning, rule-based posture correction, human-friendly feedback, and progress storage, the project attempts to create a complete AI fitness assistant instead of only a proof-of-concept activity recognition model.




\chapter{Literature Review}

\section{Introduction to Literature Review}
Recent research in exercise monitoring and human activity recognition has increasingly focused on pose-based systems because body skeleton data provides a compact and interpretable representation of movement. In many studies, raw image processing is replaced or supported by pose landmark extraction, which reduces computational cost and improves robustness for activity analysis. The literature relevant to this project spans four major areas: human pose estimation, sequence-based action recognition, repetition counting and exercise quality assessment, and AI-based fitness coaching systems. In this work, the literature base is organized around the six paper sources provided for the project, namely \cite{ai_trainer_paper}, \cite{yoga_arxiv}, \cite{paper_1374}, \cite{ijsret_1052}, \cite{icacc2021_03031}, and \cite{visapp_2024_266}.

\section{Review of Pose Estimation Based Fitness Systems}
The paper source represented by \cite{ai_trainer_paper} motivates the use of AI-enabled workout analysis systems that combine camera input, pose interpretation, and form correction. It supports the idea that a virtual trainer must go beyond simply recognizing an action and should actively observe the quality of exercise execution. This direction is closely aligned with the present project, where pose estimation is not treated as a standalone task but as the first stage of a complete exercise assistance pipeline. The practical contribution drawn from this source is the emphasis on combining exercise recognition with user guidance.

The paper source \cite{yoga_arxiv} is relevant because yoga and exercise posture analysis both depend strongly on body alignment and correct joint positioning. Research in yoga analysis often emphasizes the need for precise pose detection, stability analysis, and posture-aware feedback rather than simple class prediction. This is important for the current system because the goal is not only to say that an exercise is being performed, but also to detect whether that exercise is being performed correctly. The influence of this work appears in the project's posture-focused design, especially in the use of skeletal landmarks and angle constraints for feedback generation.

\section{Review of LSTM and BiLSTM Based Activity Recognition}
The source \cite{paper_1374} contributes to the sequence-modeling perspective of the project. Papers in this category generally show that exercise and action recognition are inherently temporal problems because a single pose snapshot is often insufficient to distinguish one movement from another. Instead, a sequence of changing body configurations must be analyzed over time. This directly supports the use of LSTM and BiLSTM models in the proposed system. The present project adopts this idea by converting frame-wise angle features into short temporal windows and using recurrent architectures to classify exercise classes from motion sequences.

The source \cite{ijsret_1052} also informs the decision to compare sequential deep learning models. Works of this type usually demonstrate that BiLSTM models can outperform simpler one-directional recurrent networks when motion context from the full sequence is important. In this project, that insight motivates the use of BiLSTM as the primary classifier and LSTM as a baseline model. The literature therefore supports the claim that temporal deep learning is more suitable for exercise classification than isolated frame-level recognition.

\section{Review of Transformer Based Sequence Recognition}
The source \cite{visapp_2024_266} is particularly useful for understanding modern computer vision and temporal sequence recognition approaches. Recent research in this area increasingly explores attention-based models because they can learn relationships between distant parts of a sequence more effectively than traditional recurrent methods in some settings. This motivates the inclusion of a Transformer comparison model in the current project. Although the project primarily uses BiLSTM for practical reasons, the Transformer architecture is included to strengthen the academic comparison and to study whether attention-based sequence processing offers benefits on landmark-derived exercise data.

\section{Review of Posture Correction and Exercise Quality Assessment}
The conference paper source \cite{icacc2021_03031} supports the posture-quality aspect of the project. Research of this kind often examines how joint-angle analysis and skeletal interpretation can be used to evaluate whether a pose or movement satisfies expected biomechanical conditions. This directly influenced the design of the posture correction engine in the present system. Instead of relying only on black-box classification, the system uses interpretable exercise-specific rules such as insufficient depth, elbow instability, arm imbalance, and incomplete lockout. The paper therefore contributes to the project's decision to combine classification with form quality assessment.

\section{Review of Repetition Counting Systems}
Across the provided sources, especially \cite{ai_trainer_paper}, \cite{paper_1374}, and \cite{ijsret_1052}, a recurring theme is that repetition counting must be based on motion cycles rather than isolated threshold checks. Several systems in the literature use phase transitions, smoothing, and temporal validation to reduce false counts caused by jitter or noisy landmarks. The current project follows the same principle. It smooths the tracked metric, identifies movement stages such as up and down, and increments the rep count only after a complete validated cycle. This makes the rep counting engine more robust for real-time use.

\section{Gaps in Existing Literature}
Although the provided paper sources offer important insights into pose estimation, activity recognition, posture quality, and intelligent workout assistance, they usually focus on only one or two parts of the complete problem. Some focus on pose classification, some on sequence modeling, and some on posture evaluation. However, few systems integrate real-time exercise classification, repetition counting, posture correction, voice guidance, database-backed workout tracking, and detailed end-of-exercise reporting into a single deployable architecture. This gap motivates the current project.

\section{Summary of Literature and Project Positioning}
The proposed FormFix AI system builds on the most useful ideas present across the supplied literature sources \cite{ai_trainer_paper,yoga_arxiv,paper_1374,ijsret_1052,icacc2021_03031,visapp_2024_266}. From pose-centered works it adopts skeleton extraction as the primary body representation. From recurrent action-recognition studies it adopts sequence learning using LSTM and BiLSTM. From quality-assessment works it adopts joint-angle based posture evaluation. From more modern sequence modeling works it includes Transformer-based comparison. However, the project extends these ideas into a practical fitness assistant that runs on consumer hardware, supports multiple workout functions, and provides both real-time and post-exercise feedback.


\chapter{Proposed System}

\section{3.1 System Overview}
The proposed system is a hybrid AI architecture composed of perception, temporal classification, movement evaluation, feedback generation, and persistence layers. The user stands in front of a webcam and performs a supported exercise. The camera captures video frames, which are passed to the pose estimation module. The pose estimator extracts full-body landmarks and returns a structured representation of the user’s body in each frame. These landmarks are then converted into meaningful angle-based features such as elbow angle, shoulder angle, knee angle, and hip angle.

The extracted frame-wise features are stored in a short temporal sequence buffer. This buffer is used by the exercise classification model, primarily a BiLSTM, to identify which exercise is being performed. Once the exercise is known, the system activates exercise-specific rep counting and posture correction logic. The rep counter tracks motion phases such as up and down, verifies full transitions, and increments the repetition count when a complete movement is detected. Simultaneously, the posture correction engine compares the observed joint angles against exercise-specific movement expectations and flags mistakes when the form deviates from expected motion patterns.

The correction output is then passed to a user-facing feedback layer. For fast real-time guidance, local posture rules are used directly. For more natural and friendly phrasing, an optional local large language model is used to convert the technical mistake information into short coach-like cues. At the end of the exercise, the system produces a detailed report containing total repetitions, perfect repetitions, corrected repetitions, frequent mistakes, and recommendations for improvement. The workout progress is finally stored in a SQL-based backend for later review.

\section{3.2 High-Level Architecture}
The complete system can be divided into the following functional layers:

\noindent\textbullet\ Input Layer  
   Webcam captures the live video stream from the user.

\noindent\textbullet\ Pose Estimation Layer  
   MediaPipe Pose extracts body landmarks and generates a skeleton representation.

\noindent\textbullet\ Feature Extraction Layer  
   The system computes relevant joint angles and derived motion features from the landmarks.

\noindent\textbullet\ Temporal Classification Layer  
   LSTM, BiLSTM, and Transformer models are used to learn motion sequences and predict exercise classes.

\noindent\textbullet\ Repetition Counting Layer  
   Exercise-specific logic tracks the movement cycle and counts completed repetitions.

\noindent\textbullet\ Posture Correction Layer  
   Angle-based rules identify form deviations and determine the most relevant mistake.

\noindent\textbullet\ Natural Coaching Layer  
   A local LLM optionally rephrases technical corrections into human-like trainer feedback.

\noindent\textbullet\ Reporting and Persistence Layer  
   Session results, progress status, completed exercises, skipped exercises, and reports are stored through backend APIs and a database.

\section{3.3 Architecture of the Models Used}

\subsection{3.3.1 MediaPipe Pose Estimation Model}
MediaPipe Pose is used as the front-end perception model. It accepts each camera frame as input and outputs a full-body landmark set. Each landmark contains spatial coordinates and a visibility score. This model is chosen because it provides high-speed full-body estimation while remaining lightweight enough for real-time execution on a MacBook Air M2. MediaPipe is not used to classify exercises directly; instead, it creates the skeletal representation on which later stages depend.

\subsection{3.3.2 Feature Representation}
The raw landmark coordinates are transformed into angle-based features because exercise correctness is more naturally described through joint mechanics than through absolute coordinate positions. For example, elbow flexion is meaningful for bicep curls and push-ups, knee angle is meaningful for squats, and shoulder angle is meaningful for lateral raises. This feature engineering stage produces a more interpretable and compact sequence representation than raw image pixels.

\subsection{3.3.3 Baseline LSTM Model}
The baseline LSTM model is a sequential neural network composed of two LSTM layers, dropout layers, a dense hidden layer, and a softmax output. The architecture is:

- LSTM(128, return_sequences=True)  
- Dropout(0.3)  
- LSTM(64)  
- Dropout(0.3)  
- Dense(64, activation="relu")  
- Dense(num_classes, activation="softmax")

This model serves as a baseline for sequence classification. It learns the temporal evolution of exercise motion from the input feature sequences. It is useful as a benchmark because LSTM is a classic and widely accepted model for time-series classification.

\subsection{3.3.4 BiLSTM Model}
The BiLSTM model extends the LSTM architecture by processing the sequence in both forward and backward directions. Its architecture is:

- Bidirectional(LSTM(128, return_sequences=True))  
- Dropout(0.3)  
- Bidirectional(LSTM(64))  
- Dropout(0.3)  
- Dense(64, activation="relu")  
- Dense(num_classes, activation="softmax")

This is the main classification model used in the project. It was selected because exercise motion contains temporal structure that benefits from richer context. BiLSTM often performs better than a single-direction LSTM on medium-sized landmark-based datasets because it captures more sequence information without becoming too computationally expensive.

\subsection{3.3.5 Transformer Model}
The Transformer comparison model uses self-attention rather than recurrence. Its architecture includes a Transformer block with multi-head attention, feed-forward layers, dropout, normalization, global average pooling, and a softmax classifier. It is useful for research comparison because attention-based models are modern and powerful for sequence learning. However, they generally require more data and tuning. For this reason, the Transformer is included to strengthen the academic value of the project and enable comparison with recurrent models, but it is not the default real-time inference model.

\subsection{3.3.6 Rep Counting Engine}
Rep counting is not performed by a separate deep learning model. Instead, it is implemented using angle transitions, smoothing, dynamic thresholds, and movement stage logic. Each supported exercise has a target metric such as average elbow angle or average shoulder angle. The rep counter uses:

- smoothing of noisy angle signals  
- observed range calibration  
- up/down stage detection  
- minimum rep gap protection  
- threshold-based cycle completion

This method is fast, interpretable, and better suited for real-time repetition tracking than a separate classifier.

\subsection{3.3.7 Posture Correction Engine}
The posture correction engine is rule-based. It receives the detected exercise and the corresponding angle features, and then checks whether the movement satisfies exercise-specific form expectations. For example:

- bicep curls are checked for elbow instability and incomplete contraction  
- push-ups are checked for hip sag and insufficient depth  
- pull-ups are checked for incomplete height and incomplete lockout  
- barbell squats are checked for insufficient depth and knee instability  
- lateral raises are checked for insufficient height, excessive height, and arm imbalance

This approach is computationally efficient and clinically interpretable compared with a black-box posture-quality classifier.

\subsection{3.3.8 Local LLM Coaching Layer}
The local LLM layer is not used as the primary posture detector. Instead, it acts as a natural-language explanation layer. It converts technical mistake codes into trainer-like voice cues and richer post-exercise reports. This design allows the system to remain real-time because the actual biomechanical decision is made locally by the rule engine. The LLM is only used to improve communication quality and end-of-exercise explanation.

\section{3.4 Why These Models Were Used}
Each model and component in the project has a specific role. MediaPipe Pose is used because real-time skeleton extraction is a prerequisite for any posture-aware exercise system. LSTM is used as a standard sequential baseline model. BiLSTM is used as the primary exercise classifier because it models temporal motion more effectively in both directions. Transformer is used as an advanced comparison model with strong sequence-learning potential. The rep counter and posture rule engine are implemented separately because repetition counting and posture correction require explicit and stable motion logic rather than only classification. Finally, the LLM layer is used because users respond better to human-like guidance than to raw technical labels.

\section{3.5 Functional Modules}
The major modules of the proposed system are:

\noindent\textbullet\ User authentication and workout selection module  
\noindent\textbullet\ Webcam capture and real-time video processing module  
\noindent\textbullet\ Skeleton detection and landmark extraction module  
\noindent\textbullet\ Feature extraction and sequence buffering module  
\noindent\textbullet\ Exercise classification module  
\noindent\textbullet\ Rep counting module  
\noindent\textbullet\ Posture correction and voice feedback module  
\noindent\textbullet\ Exercise report generation module  
\noindent\textbullet\ User progress tracking and database storage module


\chapter{Mathematical Model and Algorithms}

\section{4.1 Mathematical Representation of Landmarks}
Let the pose landmark set extracted from one frame be represented as:

P = {p1, p2, p3, ..., pn}

where each point pi contains x, y, z coordinates and a visibility measure. These landmarks correspond to major body joints and reference points required for motion analysis.

\section{4.2 Joint Angle Calculation}
Joint angle estimation is an essential part of the system because many exercises are defined by the angular relationship between connected joints. For three points A, B, and C, the angle at point B is calculated using the cosine rule:

$\theta = \cos^{-1} \left( \frac{(A-B) \cdot (C-B)}{|A-B||C-B|} \right)$

This formula is applied repeatedly to compute exercise-relevant biomechanical features such as elbow angle, knee angle, shoulder angle, and hip angle.

\section{4.3 Feature Vector Representation}
For each frame, the extracted angles are represented as:

F = [theta1, theta2, theta3, ..., thetam]

These values form the feature vector for one frame. To model exercise motion over time, a temporal sequence of features is created:

S = [F1, F2, F3, ..., Ft]

where t is the sequence length used by the classification model.

\section{4.4 Classification Function}
The classification model receives the sequence S and predicts the probability of each supported exercise class:

$\hat{Y} = \mathrm{softmax}(W h + b)$

where h is the learned sequence representation, W is the output weight matrix, and b is the output bias. The exercise label with the highest probability is selected as the predicted class.

\section{4.5 Accuracy and Performance Metrics}
The classification accuracy is calculated as:

Accuracy = (Number of Correct Predictions / Total Number of Predictions) x 100

Precision, Recall, and F1-score are defined as:

$\mathrm{Precision} = \frac{TP}{TP+FP}$
$\mathrm{Recall} = \frac{TP}{TP+FN}$
$F1 = 2 \times \frac{\mathrm{Precision} \times \mathrm{Recall}}{\mathrm{Precision}+\mathrm{Recall}}$

Top-1 accuracy measures whether the highest-probability class is correct.  
Top-3 accuracy measures whether the true class appears in the top three predicted classes.

\section{4.6 Rep Counting Logic}
For each exercise, a selected angle metric is tracked. The system smooths this metric, estimates the observed range of motion, and defines thresholds for motion stages. If the angle moves from one stage to the opposite stage and then returns while satisfying movement constraints, the repetition count is incremented by one. This prevents noisy frame fluctuations from generating false counts.

\section{4.7 Posture Evaluation Logic}
The posture engine compares observed angle patterns with predefined exercise-specific rules. If all required constraints are satisfied, the movement is labeled as good form. If one or more rules fail, the system selects the most relevant or most severe rule violation and generates a corresponding mistake code. That mistake code is then translated into corrective feedback.

\section{4.8 Algorithm 1: Pose Detection}
\noindent\textbullet\ Capture a video frame from the webcam.  
\noindent\textbullet\ Pass the frame to MediaPipe Pose.  
\noindent\textbullet\ Extract visible body landmarks.  
\noindent\textbullet\ Return the landmark set for later feature extraction.

\section{4.9 Algorithm 2: Exercise Classification}
\noindent\textbullet\ Receive a sequence of landmark-derived features.  
\noindent\textbullet\ Pass the sequence to the temporal classifier.  
\noindent\textbullet\ Obtain class probabilities.  
\noindent\textbullet\ Select the exercise with the highest probability.

\section{4.10 Algorithm 3: Repetition Counting}
\noindent\textbullet\ Choose the exercise-specific tracking angle.  
\noindent\textbullet\ Smooth the angle history.  
\noindent\textbullet\ Detect transition between stages such as up and down.  
\noindent\textbullet\ Validate full motion cycle.  
\noindent\textbullet\ Increment repetition count only if the cycle is complete.

\section{4.11 Algorithm 4: Posture Correction}
\noindent\textbullet\ Receive predicted exercise and current angle values.  
\noindent\textbullet\ Load the relevant correction rules for that exercise.  
\noindent\textbullet\ Evaluate rule failures.  
\noindent\textbullet\ Select the most relevant mistake.  
\noindent\textbullet\ Generate a correction message.  
\noindent\textbullet\ Trigger real-time voice and UI feedback.

\section{4.12 Algorithm 5: Session Report Generation}
\noindent\textbullet\ Aggregate frame-level posture statistics.  
\noindent\textbullet\ Count perfect repetitions and corrected repetitions.  
\noindent\textbullet\ Identify dominant mistake codes.  
\noindent\textbullet\ Generate a structured report.  
\noindent\textbullet\ Present the report to the user and store progress in the backend.


\chapter{Data Structures and Implementation}

\section{5.1 Dataset Sources}
The project is based on the idea of training exercise recognition on pose-informed public datasets and curated exercise clips. The datasets considered include:

- Kinetics-400 for general action recognition
- NTU RGB+D for skeleton-based activity understanding
- UCF101 for classical action recognition
- Fitness-AQA for exercise quality assessment

These datasets are useful because they provide exercise-like action classes, video samples, and in some cases skeleton-oriented activity structure. However, not all public classes map directly to practical gym exercises. Therefore, the project uses selected class filtering and exercise-specific curation rather than blindly training on all public data.

\section{5.2 Preprocessing Pipeline}
The preprocessing pipeline converts video-based data into a sequence dataset suitable for temporal models. The steps are:

\noindent\textbullet\ Read videos from class-specific folders.  
\noindent\textbullet\ Extract frames at a manageable sampling rate.  
\noindent\textbullet\ Detect skeleton landmarks using pose estimation.  
\noindent\textbullet\ Convert landmarks into angle-based features.  
\noindent\textbullet\ Form fixed-length feature windows.  
\noindent\textbullet\ Assign exercise labels.  
\noindent\textbullet\ Save processed sequences into an `.npz` dataset file.

This pipeline standardizes input across exercises and reduces noise by transforming raw video into structured motion features.

\section{5.3 Data Structures Used}
The project uses several practical data structures:

- Dictionaries for exercise configuration and posture rule templates  
- Lists and arrays for frame features and label storage  
- Deques for smoothing recent angle values  
- JSON objects for reports and UI responses  
- SQL tables for user records, workout status, and exercise progress

These choices make the system efficient, interpretable, and easy to maintain.

\section{5.4 Training Implementation}
The training implementation is written in Python using TensorFlow and Keras. The processed dataset is loaded using NumPy. The data is split into training and validation sets using `train_test_split`. Class weights are computed to reduce the impact of dataset imbalance. Two sequential models are trained directly: the BiLSTM model and the baseline LSTM model. Early stopping and learning-rate reduction callbacks are used to improve generalization and prevent unnecessary overtraining. The final models are saved in Keras format for later inference.

\section{5.5 Real-Time Inference Implementation}
During real-time use, the system captures video frames from the webcam, extracts landmarks through MediaPipe, computes current angles, and optionally updates a short sequence buffer for the classifier. The predicted exercise activates the corresponding rep counter profile and posture correction rules. The UI then shows exercise name, reps, stage, and the latest coach cue. The frontend is designed to keep the user visible while the analysis remains smooth.

\section{5.6 Frontend and Backend}
The frontend is built using React and Vite. It provides workout selection, exercise pages, camera mode, and progress display. The backend is built using FastAPI and acts as the bridge between the UI, database, and real-time AI functions. It handles:

- authentication  
- progress updates  
- exercise status  
- frame analysis requests  
- session report generation

\section{5.7 Database Design}
The project uses a relational SQL backend to store user and workout data. The key entities are:

- User  
- WorkoutDayProgress  
- WorkoutExerciseProgress  
- WorkoutExerciseStatus

These entities allow the system to track whether a user completed or skipped an exercise, how many repetitions were logged, where the user stopped, and what remains pending for future sessions.

\section{5.8 Voice and Report Generation}
Voice correction is generated through browser-based speech synthesis in the frontend. To avoid excessive repetition, the system uses cooldown logic so that corrections are spoken only when useful. The report generation pipeline aggregates exercise performance data and optionally uses a local LLM to create a more natural summary. This improves user comprehension without sacrificing the speed of the posture analysis loop.


\chapter{Results and Performance Analysis}

\section{6.1 Evaluation Methodology}
The project is evaluated using both offline model metrics and online practical behavior. Offline evaluation measures how well the classifier predicts the correct exercise label from processed pose sequences. Online evaluation considers how well the rep counter behaves, how stable the posture correction is, and how smooth the real-time user experience feels under normal webcam conditions.

\section{6.2 Training Accuracy, Validation Accuracy, and Test Accuracy}
The project code computes training and validation metrics during the training process. During model fitting, the model is evaluated on the validation split after each epoch. This gives:

- training accuracy  
- validation accuracy  
- validation loss

After training, the evaluation script computes:

- Top-1 accuracy  
- Top-3 accuracy  
- classification report  
- confusion matrix

These metrics are important because they describe not only overall accuracy but also class-wise performance and confusion behavior. The exact numeric results should be inserted into the report from the outputs produced by the training and evaluation commands used in the project environment.

\section{6.3 Classification Report}
The classification report measures precision, recall, and F1-score for each class. Precision indicates how many predicted instances of a class were correct. Recall indicates how many actual instances of that class were correctly identified. F1-score balances both precision and recall. In practical terms, high precision means the model does not mislabel other exercises as a target class too often, while high recall means it is good at detecting the target class when it actually occurs.

\section{6.4 Confusion Matrix Analysis}
The confusion matrix is useful for identifying which exercise classes are most likely to be confused. Similar upper-body movements may show overlap if their motion profiles share joint-angle patterns. For example, push-up and dip-like patterns may overlap in elbow motion, while curls and other pulling motions may show partial similarity depending on the training data. By analyzing the confusion matrix, the researcher can decide whether more data, better feature selection, or stronger temporal separation is required.

\section{6.5 Real-Time Rep Counting Analysis}
Rep counting is one of the most important practical metrics of the system. The project uses angle smoothing, dynamic thresholding, and motion-stage validation to reduce false counts. Compared with a naive threshold-based approach, the improved rep counter performs more consistently across supported exercises. Bicep curls generally show strong performance because elbow flexion is clear and easy to model. Push-ups and pull-ups require stronger control because torso alignment and visibility can introduce additional noise. Squats require deeper movement calibration and correct camera framing for stable counting.

\section{6.6 Posture Correction Analysis}
The posture correction system operates using exercise-specific rules and therefore remains interpretable. Its performance depends on the quality of skeleton detection, camera angle, and the suitability of the chosen exercise rules. Compared with purely generic feedback systems, the project offers more targeted corrections because it knows the exercise class and uses custom angle logic. However, some complex movements and occluded poses may still require more robust biomechanical modeling or richer datasets.

\section{6.7 Role of the LLM in Result Quality}
The local LLM layer improves the style and usefulness of feedback but does not replace the core posture analysis. This is an important design decision because real-time accuracy depends mainly on pose estimation and angle evaluation, not on language generation. The project therefore uses the LLM mainly for:

- human-like correction phrasing  
- trainer-style tone  
- detailed post-exercise report generation

This improves user satisfaction while preserving low latency in the real-time loop.

\section{6.8 Comparison of Models}
The LSTM baseline is useful for establishing a reference point for sequence classification. The BiLSTM typically offers better performance because it captures richer temporal context. The Transformer model is important for academic comparison and may provide stronger performance on large and well-balanced datasets, but in practical settings BiLSTM can remain the more stable and efficient choice. Therefore, the project uses BiLSTM as the main classifier while still acknowledging the value of Transformer-based sequence learning.

\section{6.9 Practical Performance on Consumer Hardware}
One of the strengths of the project is that it is designed to run on a MacBook Air M2, which is a practical consumer laptop rather than a high-end training workstation. MediaPipe Pose enables real-time skeleton extraction, and the rep counter and posture rules are lightweight enough to run continuously. The system avoids depending on heavy per-frame LLM reasoning, which would reduce smoothness. Instead, it uses local logic for fast decisions and reserves language generation for selective correction or final reporting.

\section{6.10 Sample Result Discussion}
In real usage, the system successfully detects supported exercises, counts repetitions, provides posture cues, and generates a post-exercise report. The quality of results improves when the user remains visible in the frame, uses clear movement range, and performs the supported exercises from a suitable angle. The most reliable parts of the current system are skeleton extraction, bicep curl analysis, and structured report generation. The most sensitive parts are exercise-wise rep counting under noisy conditions and form correction for complex or partially occluded motions.


\chapter{Advantages, Applications, Comparison, and Future Scope}

\section{7.1 Advantages of the Proposed System}
The proposed system has several practical and technical advantages:

\noindent\textbullet\ It integrates multiple AI tasks into one deployable solution.  
\noindent\textbullet\ It provides real-time exercise analysis using consumer-grade hardware.  
\noindent\textbullet\ It combines exercise recognition with repetition counting and posture correction.  
\noindent\textbullet\ It supports voice guidance and natural language style feedback.  
\noindent\textbullet\ It stores user progress for continuity and review.  
\noindent\textbullet\ It uses a modular architecture, making future upgrades easier.  
\noindent\textbullet\ It balances speed and interpretability through pose-based analysis.

\section{7.2 Applications}
The project can be applied in several settings:

- home fitness guidance  
- smart gym assistance  
- student and beginner training support  
- fitness monitoring dashboards  
- exercise education tools  
- lightweight movement quality systems for general wellness

With additional validation, it may also support pre-rehabilitation or supervised home exercise monitoring, though the current project is not presented as a medical device.

\section{7.3 Comparison with Existing Research Papers}
Many related papers focus on isolated tasks such as action recognition, pose estimation, or exercise quality scoring. In contrast, the proposed system combines:

- real-time skeleton extraction  
- temporal exercise classification  
- repetition counting  
- posture correction  
- voice feedback  
- progress tracking  
- detailed end-of-exercise reporting

This makes it more application-oriented than many benchmark-centered research works. While some papers may report strong classification accuracy in controlled environments, the proposed system aims for a broader and more practical fitness assistance pipeline.

\section{7.4 Limitations of the Current System}
Although the project is practical and technically strong for a student-level deployment, it has several limitations:

\noindent\textbullet\ Performance depends on camera angle and visibility.  
\noindent\textbullet\ Some exercises are harder to analyze accurately from a single webcam.  
\noindent\textbullet\ Public datasets do not always map cleanly to gym-specific movements.  
\noindent\textbullet\ Real-time posture correction still depends on well-designed rules.  
\noindent\textbullet\ Rep counting for some exercises may need further calibration.  
\noindent\textbullet\ The local LLM improves language quality but does not improve biomechanical detection itself.

\section{7.5 Future Enhancements}
Future work can improve the project significantly. Possible upgrades include:

- adding more exercise classes such as overhead press, rows, deadlifts, and bench press  
- building a cleaner custom dataset specifically for supported exercises  
- introducing user-specific calibration for height, build, and camera angle  
- improving posture evaluation through exercise quality assessment networks  
- adding mobile deployment  
- adding weekly analytics and trend dashboards  
- improving confusion reduction with richer sequence features  
- integrating multi-camera or side-view support  
- extending the database to store long-term performance history and comparative improvement reports


\chapter{Conclusion}

The project successfully demonstrates the design and implementation of an AI-based real-time fitness trainer capable of detecting exercises, counting repetitions, evaluating posture, and guiding the user through corrective feedback. By combining pose estimation, temporal deep learning, angle-based biomechanics, and user-facing feedback mechanisms, the system addresses a practical need in home and digital fitness assistance. It provides more value than systems that only recognize actions because it also understands movement quality and offers actionable corrections.

The technical contribution of the project lies in its hybrid architecture. MediaPipe Pose provides efficient skeleton extraction, sequence models provide exercise recognition, the rep counting engine enables practical workout tracking, posture rules provide interpretable correction, and the optional local LLM layer improves communication quality through more natural feedback and structured reports. This separation of responsibilities helps the system remain efficient and understandable while still offering a polished user experience.

From an application perspective, the system demonstrates that real-time AI exercise guidance can be implemented on consumer hardware with meaningful user value. It supports safer workouts, better exercise understanding, and long-term progress tracking. Although further improvements are possible in dataset quality, exercise coverage, and posture scoring sophistication, the current project establishes a strong foundation for a deployable intelligent personal training assistant.


\chapter*{Bibliography Guidance}

Use IEEE format for all references. Include:

\noindent\textbullet\ Research papers on pose estimation  
\noindent\textbullet\ Research papers on LSTM/BiLSTM/Transformer action recognition  
\noindent\textbullet\ Papers on exercise quality assessment  
\noindent\textbullet\ Dataset references for Kinetics-400, NTU RGB+D, UCF101, and Fitness-AQA  
\noindent\textbullet\ MediaPipe documentation  
\noindent\textbullet\ TensorFlow documentation  

Example format:

[1] Author Name, "Paper Title," Journal or Conference Name, vol. x, no. x, pp. xx-xx, year.  
[2] Author Name, "Paper Title," Proceedings of Conference Name, pp. xx-xx, year.  
[3] Dataset or Organization Name, "Dataset Title," year.  


\chapter*{Appendix}

Include the following in the appendix:

\noindent\textbullet\ screenshots of the website UI  
\noindent\textbullet\ camera module screenshots  
\noindent\textbullet\ architecture diagrams  
\noindent\textbullet\ confusion matrix image  
\noindent\textbullet\ model architecture diagrams  
\noindent\textbullet\ training commands  
\noindent\textbullet\ evaluation commands  
\noindent\textbullet\ sample output JSON report  
\noindent\textbullet\ database schema or table design  
\noindent\textbullet\ sample workout progress screenshots  


\section*{Suggested Tables and Figures}

\section{Suggested Figures}
\noindent\textbullet\ Overall system architecture  
\noindent\textbullet\ BiLSTM model architecture  
\noindent\textbullet\ Baseline LSTM model architecture  
\noindent\textbullet\ Transformer model architecture  
\noindent\textbullet\ Real-time correction pipeline  
\noindent\textbullet\ Frontend camera UI screenshot  
\noindent\textbullet\ Workout report UI screenshot

\section{Suggested Tables}
\noindent\textbullet\ Supported exercise list  
\noindent\textbullet\ Model comparison table  
\noindent\textbullet\ Dataset mapping table  
\noindent\textbullet\ Accuracy metrics table  
\noindent\textbullet\ Exercise-wise result table  
\noindent\textbullet\ Database table summary


\section*{Accuracy and Metrics Section Guidance}

Insert your actual results from:

```bash
python training/train_lstm.py
python training/evaluate.py --model models/exercise_bilstm.keras
```

Metrics to include:

- Training Accuracy
- Validation Accuracy
- Test Accuracy / Top-1 Accuracy
- Top-3 Accuracy
- Precision
- Recall
- F1-score
- Confusion Matrix

Explain them as:

Accuracy = correct predictions divided by total predictions.  
Precision = reliability of positive predictions.  
Recall = ability to identify actual positives.  
F1-score = balance between precision and recall.


\section*{Short Ready-to-Use Contents Page}

Chapter 1: Introduction  
Chapter 2: Literature Review  
Chapter 3: Proposed System  
Chapter 4: Mathematical Model and Algorithms  
Chapter 5: Data Structures and Implementation  
Chapter 6: Results and Performance Analysis  
Chapter 7: Advantages, Applications, Comparison and Future Scope  
Chapter 8: Conclusion  
Bibliography  
Appendix


\begin{thebibliography}{9}
\bibitem{ai_trainer_paper} \textit{AI\_Trainer\_ResearchPaper.pdf}, project reference paper provided by the author.
\bibitem{yoga_arxiv} \textit{Yoga\_Arxiv.pdf}, project reference paper provided by the author.
\bibitem{paper_1374} \textit{1374-1381-1.pdf}, project reference paper provided by the author.
\bibitem{ijsret_1052} \textit{IJSRET\_V11\_issue3\_1052.pdf}, project reference paper provided by the author.
\bibitem{icacc2021_03031} \textit{itmconf\_icacc2021\_03031.pdf}, project reference paper provided by the author.
\bibitem{visapp_2024_266} \textit{VISAPP\_2024\_266\_CR.pdf}, project reference paper provided by the author.
\end{thebibliography}

\end{document}
